\documentclass[11pt,a4paper]{article}
\usepackage[margin=25mm]{geometry}
\usepackage[T1]{fontenc}
\usepackage{lmodern,amsmath,amssymb,amsthm,enumitem,microtype}
\usepackage[colorlinks=true,urlcolor=blue,citecolor=blue,linkcolor=blue]{hyperref}
\usepackage{fancyhdr}
\pagestyle{fancy}
\fancyhf{}
\lhead{Two-sided 60/40 majority}
\rhead{21 October 2026}
\cfoot{\thepage}
\setlength{\headheight}{14pt}
\setlength{\parskip}{3pt}
\setlength{\emergencystretch}{2em}
\newtheorem{theorem}{Theorem}
\newtheorem{lemma}[theorem]{Lemma}
\newtheorem{corollary}[theorem]{Corollary}
\newtheorem{proposition}[theorem]{Proposition}
\theoremstyle{remark}
\newtheorem{remark}[theorem]{Remark}
\newcommand{\sgnplus}{\operatorname{sgn}_{+}}
\newcommand{\err}{\operatorname{err}}
\newcommand{\unl}{\operatorname{unl}}
\newcommand{\norm}[1]{\left\lVert#1\right\rVert}
\newcommand{\ip}[2]{\langle #1,#2\rangle}
\newcommand{\1}{\mathbf 1}
\newcommand{\RP}{\mathrm{RP}}
\newcommand{\NP}{\mathrm{NP}}
\title{Two-sided 60/40 majority:\\prediction, counting, and compatibility}
\author{Grisha Pochuev\\\small Independent Researcher\\
\small\href{https://orcid.org/0009-0005-0078-0857}{ORCID: 0009-0005-0078-0857}\\
\small\href{mailto:n_854@mail.ru}{n\_854@mail.ru}}
\date{21 October 2026\\[4pt]\small Manuscript, version 1.0}
\hypersetup{pdftitle={Two-sided 60/40 majority: prediction, counting, and compatibility},
pdfauthor={Grisha Pochuev},pdfsubject={Explicit learning guarantees and compatibility hardness},
pdfkeywords={PAC learning, majority with margins, perceptron, sample compression, NP-completeness}}

\begin{document}
\maketitle

\begin{abstract}
We give self-contained prediction guarantees for two precise interpretations of
two-sided majority with a 60/40 gap. When the percentages are taken within the
positive and negative indicator groups, centering the Boolean coordinates gives
a Euclidean margin of $1/5$. A perceptron with an elementary stopping rule learns
from at most $26\lceil\varepsilon^{-1}\ln(26/\delta)\rceil$ labeled examples,
independently of dimension. When the percentages are taken among the active
indicators, we prove the distribution-free counting inequality
$\ln K_\varepsilon\ge \varepsilon n/5-\ln2$, where $K_\varepsilon$ is the number
of hypotheses whose incompatibility is at most $\varepsilon$. A normalized
perceptron consequently uses $O(B^2\log(2+B))$ labels, where
$B=(\ln K_\varepsilon+\ln(1/\delta))/\varepsilon$ is the raw-count benchmark.
Both procedures have polynomial bit complexity and require no additional
unlabeled observations. The first appendix gives a polynomial-time learner
returning a signed-literal majority, with a sample bound proved by compressed
descriptions. The second proves NP-completeness of finding an exactly compatible
partition and derives an $\NP=\RP$ consequence for general polynomial-time
learning with an output-compatibility guarantee. These are distinct mathematical
requirements. The classical learning mechanisms are identified explicitly,
and the results are compared with the polynomial prediction criterion in the
Blum--Balcan question under stated normalization and output conventions.
\end{abstract}

\section{The learning criterion and the two promises}

Write $[n]=\{1,\ldots,n\}$, and set $\sgnplus(t)=1$ when $t\ge0$ and $-1$
otherwise. Throughout, $\mathcal D$ is a distribution on labeled examples
$(x,y)$, $D$ is its marginal on $x$, and
\[
 \err_{\mathcal D}(h)=\Pr_{(X,Y)\sim\mathcal D}[h(X)\ne Y].
\]
All examples used for learning are independent and identically distributed
(iid). We take $0<\varepsilon<1$ and $0<\delta\le1/2$ unless stated otherwise.
There is no independence assumption between coordinates, and no balance
assumption on the class probabilities or the unknown group sizes.
An efficient learner runs in time polynomial in $n$, $1/\varepsilon$,
$\log(1/\delta)$, and the bit lengths of rational parameter encodings.
Calls to a fresh iid example source are counted as labeled observations;
reading a returned $n$-bit point costs $O(n)$ operations. A prediction guarantee
alone does not require the returned classifier to belong to the target class
or its defining partition to remain compatible with $D$.

For a finite hypothesis class $C$ and a compatibility function
$\chi:C\times X\to[0,1]$, define
\[
 \unl_D(h)=\mathbb E_{X\sim D}[1-\chi(h,X)],
 \qquad
 C_{D,\chi}(\varepsilon)=\{h\in C:\unl_D(h)\le\varepsilon\}.
\]
The raw-count labeled-sample benchmark in \cite[Section 3]{bb2007} and
\cite[Theorem 4]{bb2010} is, up to absolute constants,
\begin{equation}\label{eq:benchmark}
 B(C,D,\chi;\varepsilon,\delta)
 =\frac{\ln|C_{D,\chi}(\varepsilon)|+\ln(1/\delta)}{\varepsilon}.
\end{equation}
Replacing $\ln(1/\delta)$ by $\ln(2/\delta)$ changes the comparison only by
absolute constants in the stated range. The 2007 question explicitly allows
an efficient algorithm using this many labels or a polynomial in this many.
Its footnote proposes a cover-based bound as a desirable refinement; it does
not make that refinement the criterion in \eqref{eq:benchmark}.

We distinguish the following two formal promises.

\paragraph{Group normalization.}
The coordinates are partitioned as $P\sqcup N=[n]$. Put $a=|P|$, $b=|N|$, and
\[
 p(x)=\sum_{i\in P}x_i,\qquad q(x)=\sum_{i\in N}x_i.
\]
For $x\in\{0,1\}^n$, the labeled distribution satisfies, almost surely,
\begin{align}
 y=1&:\quad p(x)\ge3a/5,\qquad q(x)\le2b/5, \label{eq:groupplus}\\
 y=-1&:\quad p(x)\le2a/5,\qquad q(x)\ge3b/5. \label{eq:groupminus}
\end{align}
Here the percentages are relative to the complete groups, including inactive
coordinates. The learner knows neither the partition nor its sizes.

\paragraph{Active normalization.}
For $x\in\{0,1\}^n\setminus\{0\}$, let $k(x)=\norm{x}_1$. There is an unknown
$s^*\in\{-1,1\}^n$ such that, almost surely,
\begin{equation}\label{eq:activepromise}
 y\ip{s^*}{x}\ge k(x)/5.
\end{equation}
Thus at least 60\% of the active indicators have the sign of the label.
The associated class and compatibility function are
\begin{equation}\label{eq:activeclass}
 h_s(x)=\sgnplus(\ip{s}{x}),\qquad
 \chi_{\rm act}(s,x)=
 \1\{|\ip{s}{x}|\ge k(x)/5\}.
\end{equation}
This is the restriction to nonzero examples of the explicit majority model
in \cite[Appendix A.2]{bb2014supp}. The exclusion is stated here because the
normalized update divides by $k(x)$.

These are different promises. A positive example with one active positive
indicator satisfies \eqref{eq:activepromise}, but fails
\eqref{eq:groupplus} when $|P|\ge2$. Conversely, with
$(a,b,p,q)=(10,90,6,36)$, the group promise holds for a positive example, while
the active signed sum is $p-q=-30$. We prove a prediction theorem for each
promise rather than identify the two.

\paragraph{Organization and scope.}
Section 3 specializes the classical Euclidean perceptron bound to the group
promise. Section 4 compares the active promise with its full-function
approximate-compatibility count. Appendix A supplies an explicit sign-vector
construction and compressed-description bound. Appendix B proves exact
compatibility hardness. All probability guarantees refer to iid sampling and
one fixed realizable target. None asserts recovery of the unknown partition,
robustness to arbitrary label noise, or a count bound with a separately
vanishing compatibility tolerance. The historical question is discussed in
Section 5; no award acceptance or exhaustive priority claim is inferred.

\section{Two elementary learning lemmas}

The following proofs supply both the optimization and statistical ingredients.
They also specify a finite stopping rule rather than merely an expected-error
conversion of an online algorithm.
The first is the classical perceptron analysis associated with
Novikoff~\cite{novikoff1962}. The consecutive-correct conversion in the second
is already used by Balcan et al.~\cite[Section 3.2]{bbel2013}; we include a
proof to fix the adaptive stopping convention and constants, not to claim
that conversion as new.

\begin{lemma}[Perceptron update bound]\label{lem:perceptron}
Suppose every feature-label pair $(v,y)$ in a sequence satisfies
\[
 \norm{v}_2\le R,\qquad y\ip{u}{v}\ge\rho>0
\]
for one fixed vector $u$ with $\norm{u}_2\le U$. Initialize $w=0$, predict
$\sgnplus(\ip{w}{v})$, and on an incorrect prediction update $w\leftarrow w+yv$.
There are at most $(UR/\rho)^2$ updates.
\end{lemma}
\begin{proof}
Index only the updates by $j=1,\ldots,M$. On a mistake,
$y_j\ip{w_{j-1}}{v_j}\le0$, including the specified tie convention. Consequently,
\[
 \norm{w_j}_2^2
 =\norm{w_{j-1}}_2^2+2y_j\ip{w_{j-1}}{v_j}+\norm{v_j}_2^2
 \le\norm{w_{j-1}}_2^2+R^2.
\]
Hence $\norm{w_M}_2^2\le MR^2$, whereas
\[
 \ip{u}{w_M}=\sum_{j=1}^M y_j\ip{u}{v_j}\ge M\rho.
\]
Cauchy--Schwarz gives $M\rho\le UR\sqrt M$. For $M>0$, squaring after division
by $\sqrt M$ proves the assertion; $M=0$ is immediate.
\end{proof}

\begin{lemma}[Consecutive-correct stopping rule]\label{lem:streak}
Consider an online learner which changes its classifier only on an incorrect
prediction and which makes at most $M_0$ mistakes on every promised sequence.
Let
\[
 L=\left\lceil\frac{\ln((M_0+1)/\delta)}{\varepsilon}\right\rceil.
\]
Run it on fresh iid labeled examples, resetting a counter to zero on every
mistake and incrementing the counter on every correct prediction. Stop when
the counter reaches $L$. The procedure uses at most $(M_0+1)L$ examples and
returns a classifier of error at most $\varepsilon$ with probability at least
$1-\delta$.
\end{lemma}
\begin{proof}
A phase is an interval between counter resets, beginning initially or
immediately after a mistake. The classifier is fixed within the phase,
although a mistake need not change its values everywhere on the domain.
Each failed phase contains
at most $L-1$ correct predictions followed by one mistake. The final phase
contains $L$ correct predictions. There are at most $M_0$ failed phases and
one final phase, proving the deterministic sample bound and termination.

Condition on the complete history at the beginning of a reached phase. Its
classifier $h$ is then fixed and all subsequent observations are independent
of that history. If $\err_{\mathcal D}(h)>\varepsilon$, the conditional
probability of ending this phase with $L$ correct predictions is at most
\[
 (1-\err_{\mathcal D}(h))^L
 \le (1-\varepsilon)^L
 \le e^{-\varepsilon L}
 \le\frac{\delta}{M_0+1}.
\]
Average over histories and take a union bound over the at most $M_0+1$ possible
phase indices. The probability of accepting any classifier of error greater
than $\varepsilon$ is at most $\delta$.
\end{proof}

A fixed-sample implementation may draw $(M_0+1)L$ iid examples in advance and
process the necessary prefix. The implementation achieves this same upper
bound without a real logarithm oracle. Write $H=M_0+1$, and define
\begin{equation}\label{eq:exactstreak}
 L_* = \min\{\ell\in\mathbb Z_{\ge1}:H(1-\varepsilon)^\ell\le\delta\}.
\end{equation}
The preceding proof works directly with $L_*$, and
$L_*\le\lceil\ln(H/\delta)/\varepsilon\rceil$.
For $\varepsilon=p/q$ and $\delta=r/s$ with positive integer numerators and
denominators, its test is exactly
\[
 Hs(q-p)^\ell\le r q^\ell.
\]
Set $j=\lceil\log_2(H/\delta)\rceil$ by integer comparisons and
$U=\lceil j/\varepsilon\rceil$. Then $U$ is a sufficient upper bound for
$L_*$. Binary search on $[0,U]$ with the displayed monotone integer test
finds $L_*$ exactly. The intermediate integers use
$O(U\log(q+1)+\log(Hs+r))$ bits, which is polynomial in the stated input
parameters. The shared implementation is \path{pac_thresholds.py}.
Explicitly capping observations and updates gives a polynomial runtime also
on inputs outside the promise.

\section{Group normalization: a dimension-independent predictor}

For analysis, define the target signs and centered features by
\[
 s_i^*=\begin{cases}1&i\in P,\\-1&i\in N,\end{cases}
 \qquad z(x)=2x-\1.
\]
Both $s^*$ and $z(x)$ have Euclidean norm $\sqrt n$, and
\begin{equation}\label{eq:centeridentity}
 \ip{s^*}{z(x)}
 =(2p(x)-a)-(2q(x)-b).
\end{equation}
For a positive example, the first parenthesis is at least $a/5$ and the
second is at most $-b/5$. For a negative example their inequalities reverse.
Thus
\begin{equation}\label{eq:groupmargin}
 y\ip{s^*}{z(x)}\ge n/5
\end{equation}
almost surely. Equivalently, $s^*/\sqrt n$ and $z(x)/\sqrt n$ are unit
vectors with signed inner product at least $1/5$.

For $n\ge1$, the two branches \eqref{eq:groupplus}--\eqref{eq:groupminus}
cannot both hold. Indeed, each nonempty group would require its active count
to be both at least three fifths and at most two fifths of its size.
Thus the promise determines labels on its support. If one group is empty,
\eqref{eq:groupmargin} and the argument below are unchanged.

\begin{theorem}[Group-normalized prediction]\label{thm:group}
Under \eqref{eq:groupplus}--\eqref{eq:groupminus}, there is a learner returning
\[
 \widehat h_w(x)=\sgnplus\!\left(\sum_{i=1}^n w_i(2x_i-1)\right),
 \qquad w\in\mathbb Z^n,
\]
with population error at most $\varepsilon$ with probability at least
$1-\delta$. It uses no additional unlabeled examples and at most
\begin{equation}\label{eq:groupsample}
 m_{\rm grp}
 =26\left\lceil\frac{\ln(26/\delta)}{\varepsilon}\right\rceil
\end{equation}
labeled examples. It requires $O(nm_{\rm grp})$ elementary arithmetic operations,
and every returned coefficient satisfies $|w_i|\le25$.
\end{theorem}
\begin{proof}
Use the perceptron on $z(x)$, with the integral update
$w\leftarrow w+yz(x)$ on each prediction mistake. Apply
Lemma~\ref{lem:perceptron} with $U=R=\sqrt n$ and $\rho=n/5$.
The bound is exactly 25 updates. Lemma~\ref{lem:streak} with $M_0=25$
gives \eqref{eq:groupsample} and the confidence guarantee.
Each coordinate changes by $1$ or $-1$ on an update, hence $|w_i|\le25$.
\end{proof}

\paragraph{Why centering is necessary.}
In the example $(a,b,p,q)=(10,90,6,36)$, the uncentered score is $-30$, but
\[
 2(p-q)-(a-b)=2(6-36)-(10-90)=20=n/5.
\]
The unknown offset is represented automatically through the centered features.
No estimation of $a$ or $b$ is required.

\paragraph{Bit complexity.}
Each score has absolute value at most $25n$ and requires
$O(\log(n+1))$ bits. The counter needs $O(\log L)$ bits. Thus the exact
integer implementation has polynomial bit complexity; a straightforward bound
is $O(nm_{\rm grp}\log((n+1)(m_{\rm grp}+1)))$ bit operations, in addition to the polynomial
work of computing the stopping threshold from the rational input parameters.

\paragraph{Comparison with the raw-count benchmark.}
Whenever a perfectly compatible target belongs to $C$,
$|C_{D,\chi}(\varepsilon)|\ge1$. In the range $\delta\le1/2$,
\[
 \frac{1+\ln(1/\delta)}{\varepsilon}
 =O\!\left(B(C,D,\chi;\varepsilon,\delta)\right).
\]
Therefore \eqref{eq:groupsample} is $O(B)$, and in particular meets the
polynomial allowance in the 2007 criterion. This comparison needs neither
an estimate of the number of compatible partitions nor a specified extension
of the target classifier away from the compatible support. The probability
distribution with its given labels is the premise of the theorem.

\section{Active normalization: a counting lemma and a predictor}

In this section use precisely the class and compatibility function
\eqref{eq:activeclass}. Set
\[
 K_\varepsilon
 =\big|\{s\in\{-1,1\}^n:\unl_D(s)\le\varepsilon\}\big|,
 \qquad
 B_{\rm act}=\frac{\ln K_\varepsilon+\ln(1/\delta)}{\varepsilon}.
\]
Different sign vectors give different functions on the full instance space:
$h_s(e_i)=s_i$. Thus $K_\varepsilon$ counts distinct functions, not merely
different descriptions of one function.

\begin{lemma}[Universal lower bound on the raw count]\label{lem:count}
For every distribution $D$ on $\{0,1\}^n\setminus\{0\}$ and every
$0<\varepsilon<1$,
\begin{equation}\label{eq:countbound}
 K_\varepsilon\ge
 \frac{1}{2(1-\varepsilon/5)^n},
 \qquad
 \ln K_\varepsilon\ge\frac{\varepsilon n}{5}-\ln2.
\end{equation}
In particular,
\begin{equation}\label{eq:dimension}
 n\le\frac{5(\ln K_\varepsilon+\ln2)}{\varepsilon}.
\end{equation}
\end{lemma}
\begin{proof}
Choose a random sign vector $s$ with independent coordinates and
$\Pr[s_i=-1]=p=\varepsilon/5$. For a fixed nonempty active set
$S=\{i:x_i=1\}$, let $R_S$ be the number of minus signs in $S$.
Incompatibility requires
\[
 \big||S|-2R_S\big|<|S|/5,
\]
and therefore requires $R_S>2|S|/5$. Since $\mathbb E R_S=p|S|$, Markov's
inequality gives
\[
 \Pr_s[\chi_{\rm act}(s,x)=0]
 \le\Pr_s[R_S>2|S|/5]\le\frac{5p}{2}=\frac{\varepsilon}{2}.
\]
Average over $x\sim D$ to obtain
$\mathbb E_s\unl_D(s)\le\varepsilon/2$. Applying Markov's inequality again
shows that the set $\unl_D(s)\le\varepsilon$ has probability at least $1/2$
under this product measure.

Here $p<1/2$, so every individual sign vector has probability at most
$(1-p)^n$: an atom with $r$ minus signs has probability
$p^r(1-p)^{n-r}\le(1-p)^n$. It follows that
$K_\varepsilon(1-p)^n\ge1/2$. Taking logarithms and using
$-\ln(1-p)\ge p$ proves both assertions.
\end{proof}

There is also a useful discrete form. If
$r=\lfloor\varepsilon n/5\rfloor$, choose the minus coordinates uniformly
among the $r$-subsets of $[n]$. The same two Markov arguments apply, because
$\mathbb E R_S=r|S|/n\le\varepsilon|S|/5$. At least half these subsets are
approximately compatible. Consequently,
\begin{equation}\label{eq:binomialcount}
 K_\varepsilon\ge\tfrac12\binom nr.
\end{equation}

\begin{theorem}[Active-normalized prediction]\label{thm:active}
Under \eqref{eq:activepromise}, a learner returning a rational-weight linear
classifier uses no additional unlabeled observations and at most
\begin{equation}\label{eq:activesample}
 m_{\rm act}
 =(25n+1)\left\lceil
 \frac{\ln((25n+1)/\delta)}{\varepsilon}\right\rceil
\end{equation}
labeled examples. With probability at least $1-\delta$, the returned classifier
has error at most $\varepsilon$. The learner has polynomial bit complexity, and
\begin{equation}\label{eq:activeB}
 m_{\rm act}=O\!\left(B_{\rm act}^2\log(2+B_{\rm act})\right).
\end{equation}
\end{theorem}
\begin{proof}
Use the normalized feature
\[
 v(x)=\frac{x}{k(x)}.
\]
It satisfies $\norm{v(x)}_2=1/\sqrt{k(x)}\le1$ and, by the promise,
$y\ip{s^*}{v(x)}\ge1/5$. Since $\norm{s^*}_2=\sqrt n$,
Lemma~\ref{lem:perceptron} gives at most $25n$ mistakes for the update
$w\leftarrow w+yv(x)$. Prediction may use
$\sgnplus(\ip{w}{x})$, since $k(x)>0$.
Lemma~\ref{lem:streak} proves \eqref{eq:activesample} and the error guarantee.

Both constant sign vectors are perfectly compatible, so $K_\varepsilon\ge2$.
Lemma~\ref{lem:count} now implies
\[
 n\le\frac{10\ln K_\varepsilon}{\varepsilon}\le10B_{\rm act},
 \qquad
 \frac1\varepsilon\le\frac{B_{\rm act}}{\ln2},
 \qquad
 \frac{\ln(1/\delta)}{\varepsilon}\le B_{\rm act}.
\]
Substitute these inequalities into \eqref{eq:activesample}. More explicitly,
\[
 m_{\rm act}\le
 (250B_{\rm act}+1)
 \left[
 \frac{B_{\rm act}}{\ln2}\ln(250B_{\rm act}+1)
 +B_{\rm act}+1
 \right].
\]
This is \eqref{eq:activeB}.
\end{proof}

\paragraph{Exact implementation.}
Let
\[
 Q=\operatorname{lcm}(1,\ldots,n)
\]
and store the integer vector $W=Qw$. An update becomes
\[
 W\longleftarrow W+y\,\frac{Q}{k(x)}x,
\]
and prediction uses $\sgnplus(\ip{W}{x})$. There are at most $25n$ updates,
so $|W_i|\le25nQ$. The estimate $Q\le n!$ gives
$\log Q=O(n\log(n+1))$. Computing $Q$ by integer gcd and multiplication,
and all subsequent predictions and updates, therefore take polynomially many
bit operations. The integral stopping threshold following
Lemma~\ref{lem:streak} avoids a logarithm oracle. No optimization or counting
oracle is used.

\paragraph{Nonempty indicator groups.}
If both groups must be nonempty, let $K_\varepsilon^{\rm ne}$ count only the
nonconstant sign vectors. Under the realizability assumption, $s^*$ and
$-s^*$ are distinct admissible vectors, so $K_\varepsilon^{\rm ne}\ge2$.
The full count is exactly $K_\varepsilon^{\rm ne}+2$, because the two removed
constant vectors are always compatible. Equation~\eqref{eq:dimension} gives
\[
 n\le
 \frac{5\ln(2(K_\varepsilon^{\rm ne}+2))}{\varepsilon}
 \le\frac{15\ln K_\varepsilon^{\rm ne}}{\varepsilon}.
\]
For the second inequality, use
$2(K+2)\le4K\le K^3$ for $K\ge2$.
The proof of \eqref{eq:activeB} therefore remains valid with the nonempty-group
count and an altered absolute constant.

\paragraph{What the counting argument uses.}
The same $\varepsilon$ is used for allowed incompatibility and desired
prediction error. The proof does not replace $K_\varepsilon$ by $K_0$, by
a number of equivalence classes on the support, or by a covering number.
For example, if $D$ is uniform on the pairs
$x=e_i+e_{i+1}$ of a connected path, perfect compatibility requires adjacent
signs to agree, so $K_0=2$ for arbitrarily large $n$. No analogue of
\eqref{eq:dimension} with a fixed zero incompatibility tolerance follows.
Nor does the argument provide a bound solely in terms of a benchmark with an
independent tolerance $\eta$ approaching zero. The theorem proves the
raw-count, same-$\varepsilon$, polynomial criterion as stated; an $O(B_{\rm act})$
bound or a cover-based refinement is a stronger question.

\section{Prediction, representation, and compatibility}

The preceding results settle the two explicit prediction formulations in this
note under the raw-count criterion allowing a polynomial in the benchmark.
The group theorem is stronger in sample size: it is $O(B)$ rather than merely
polynomial in $B$. Both procedures are standard perceptron updates specialized
to the promises; their reduction, stopping rule, and count comparison are
proved here.

There are three different output requirements:
\begin{enumerate}[itemsep=2pt]
\item a classifier with small prediction error;
\item a classifier represented by a sign vector with a specified decision rule;
\item an explicit partition with small incompatibility under the original 60/40
inequalities.
\end{enumerate}
Appendix~\ref{app:proper} proves the second requirement for the signed-literal
rule associated with centered coordinates. Appendix~\ref{app:hardness}
shows why the third requirement has a substantially different computational
status: exact finite consistency is NP-complete, and a general polynomial-time
learner guaranteeing small output incompatibility would imply $\NP=\RP$.

\paragraph{Relation to the published wording.}
The majority paragraph in \cite[Section 3.2]{bb2007}, repeated in
\cite[Section 6.2]{bb2010}, describes the percentages in words. It does not give
a formula specifying the majority classifier on incompatible examples. Thus the
centered signed-literal extension used here is explicitly our defined decision
rule, not a formula attributed to that paragraph. On every example satisfying
the group promise, it agrees with the prescribed label. For nonempty groups it
also agrees there with positive-group majority and inverted negative-group
majority. Outside that compatible region, these extensions can differ.

The 2014 paper \cite{bb2014} describes its majority setting as a modified form,
and its supplement \cite[Appendix A.2]{bb2014supp} explicitly uses fractions
of the active indicators. Our second theorem analyzes that precise class.

\paragraph{Relation to closest-string approximation and prior methods.}
For a labeled sample, put $u_j=y_j(2x_j-\1)\in\{-1,1\}^n$.
The group promise implies
\[
 d_H(s^*,u_j)=\frac{n-\ip{s^*}{u_j}}2\le\frac{2n}{5}.
\]
The closest-string approximation scheme of Li, Ma, and
Wang~\cite[Theorem 2.6]{lmw2002}, with approximation ratio $9/8$, therefore
returns a sign vector $s$ with $d_H(s,u_j)\le9n/20$ and hence
$y_j\ip{s}{2x_j-\1}\ge n/10$ on every training observation.
This is our direct application of their theorem, not a majority-learning
claim made in that paper. It already gives polynomial-time fitting with a
constant margin. Combining it with uniform margin bounds, such as
\cite[Theorem 3]{bb2014} for the Euclidean unit ball, yields dimension-independent
PAC learning as well. Accordingly, no priority claim for the existence of a
sign-vector learner is made here. Appendix A supplies a different finite-bit
construction and an explicit compressed-description bound, with all constants
proved. The general use of sample compression for learnability is classical;
see Floyd and Warmuth~\cite{fw1995}.

\paragraph{Related compatibility hardness.}
Balcan et al.~\cite[Supplement, Section 2, Theorem 4]{bbel2013supp} already
prove that finding a consistent and compatible two-sided disjunction is
NP-hard. That problem has a different concept class and compatibility
condition. Appendix B gives a self-contained reduction for the fixed group
60/40 inequalities, including balanced points and a connected Hamming-4 graph;
it does not claim the first hardness result for compatibility-constrained
learning in general.

\paragraph{Learning outside the target class.}
The framework in \cite[Section 2]{bb2010} states a prediction-error goal.
Its finite-class sufficient method uses an empirically correct and compatible
member of $C$; that sufficient method is not itself a restriction requiring
every learner to output such a member. As further evidence about the model,
\cite[Section 3.2, Algorithm 2 and Theorem 2]{bbel2013} explicitly studies a
non-proper learner that returns a low-error hypothesis not necessarily
consistent or compatible. That later result has its own distributional
parameters; we cite it for the nature of its output guarantee, not as evidence
that another prize was awarded or that all versions of the 2007 problems were
already resolved.

The mathematical conclusions are for the explicit promises and resource
criterion stated here. They do not identify a full-function count with a
distributional covering number, or the label-error tolerance with an
independently specified smaller compatibility tolerance. Priority beyond the
comparisons above, any intended stronger interpretation of the historical
question, and acceptance of a submission are not established by these proofs.

\section*{Reproducibility and assistance}

The accompanying package places Python files in \path{code/}, complete
verification records in \path{results/}, and internal review reports in
\path{audit/}. The weighted predictors are implemented in \path{majority_6040.py} for group
normalization and \path{majority_active.py} for active normalization.
The representation learner is \path{proper_signed_majority.py}.
Finite checks are supplied by \path{verify_6040.py} for the first predictor,
\path{verify_active.py} for the active count and predictor,
\path{check_proper.py} for the sign-vector construction, and
\path{verify_reduction.py} for the compatibility reduction.
The shared exact stopping threshold and its minimality are checked by
\path{verify_thresholds.py}. The package-level \path{run_all_checks.py}
runs these checks in an isolated temporary directory and records their outputs.

The initial exhaustive margin checks, through $n=10$, examine
1,398,100 partition--example pairs, including 406,868 admissible labeled
pairs, and attain the claimed minimum normalized margin $1/5$.
They include empty and unequal groups. The active-version checks include
exact rational finite counts, stopping-probability inequalities, and 64
iid learning trials on explicitly specified finite distributions.
Supplemental checks cover all clipped coefficient values and five-bit rounding
inputs, small-domain enumeration, large-dimensional correction and rounding,
and three finite instances of the compatibility reduction. These checks support
the implementations and expose edge cases; no universal claim in this note is
inferred from finite experiments or a solver status. The proofs below and above
supply the general guarantees.

This manuscript and its code were prepared with assistance from ChatGPT
(GPT-6 Astra), including separately assigned mathematical and source-review
passes. These are internal AI-assisted checks, not external peer review or
formal verification by a proof assistant. All claims are supported by the
mathematical arguments in the manuscript, rather than by the review process.

\appendix
\section{Returning a signed-literal majority}\label{app:proper}

This appendix gives a stronger representation guarantee for the group-normalized
prediction task. Its decision rule is specified throughout; output compatibility
is not asserted.

\begin{theorem}[A sign-vector output with a compressed-description proof]
\label{thm:proper}
Suppose $n\ge1$, $0<\varepsilon<1$, $0<\delta\le1/2$, and there is
$s^*\in\{-1,1\}^n$ with
\[
 y\ip{s^*}{2x-\1}\ge n/5
\]
almost surely. Put
\begin{align}
 q&=\left\lceil\log_2(2/\delta)\right\rceil, \label{eq:properq}\\
 r&=\min\{j\in\mathbb Z_{\ge0}:
       \varepsilon 2^j\ge512(j+2)+q\},\qquad m=2^r. \label{eq:properm}
\end{align}
There is a randomized learner using $m$ iid labeled observations, no additional
unlabeled observations, and at most $5nq$ independent random bits, which runs
in time polynomial in $n,1/\varepsilon,\log(1/\delta)$ and returns
$s\in\{-1,1\}^n$ such that
\[
 h_s^{\rm lit}(x)=\sgnplus(\ip{s}{2x-\1})
\]
has error at most $\varepsilon$ with probability at least $1-\delta$.
Moreover,
\begin{equation}\label{eq:properexplicit}
 m\le
 \frac{2048\big(q+\lceil\log_2(1/\varepsilon)\rceil+12\big)}{\varepsilon}.
\end{equation}
\end{theorem}

The returned vector defines $P=\{i:s_i=1\}$ and $N=\{i:s_i=-1\}$.
The rule is the majority of the $n$ literals $x_i$ for $i\in P$ and
$1-x_i$ for $i\in N$:
\begin{equation}\label{eq:literalrule}
 h_s^{\rm lit}(x)=1
 \quad\Longleftrightarrow\quad
 p(x)+|N|-q(x)\ge n/2.
\end{equation}
This rule may differ from majority over $P$ alone on incompatible examples.
The premise of Theorem~\ref{thm:proper} follows from the original group
promise by \eqref{eq:centeridentity}--\eqref{eq:groupmargin}.

\subsection{An exact-arithmetic construction}

Before drawing the sample, independently draw $q$ tables
$R^{(1)},\ldots,R^{(q)}$, each consisting of $n$ independent uniform integers
from $\{0,\ldots,31\}$. Each entry uses exactly five fair bits. This choice is
independent of the labeled observations. The same distribution of tables can
be implemented by lazy access to independent prespecified bits.

Draw $(x_j,y_j)$ for $1\le j\le m$, set $z_j=2x_j-\1$, and put
$b=\lceil\log_2(2m)\rceil$.

\paragraph{Small-dimensional branch: $n<512b$.}
Enumerate the sign vectors in a fixed order and return the first satisfying
\begin{equation}\label{eq:properenum}
 5y_j\ip{s}{z_j}\ge n\qquad(1\le j\le m).
\end{equation}
The promise supplies such a vector. On an input outside the promise with no
such vector, return an arbitrary vector with a failure flag.

\paragraph{Large-dimensional branch: $n\ge512b$.}
Initialize $a=0\in\mathbb Z^n$. If some sample observation satisfies
$y_j\ip{a}{z_j}<2n$, perform the coordinatewise correction
\begin{equation}\label{eq:clip}
 a_i\longleftarrow
 \max\{-16,\min\{16,a_i+y_jz_{j,i}\}\}\qquad(1\le i\le n),
\end{equation}
and restart the scan. Stop correcting when there is no violation. Impose a
cap of 183 corrections, returning failure if the cap is reached;
this only specifies bounded behavior outside the promise.

For each table $\ell=1,\ldots,q$, form a sign vector by
\begin{equation}\label{eq:round}
 s_i^{(\ell)}
 =\begin{cases}
 1 & R_i^{(\ell)}<a_i+16,\\
 -1 & R_i^{(\ell)}\ge a_i+16.
 \end{cases}
\end{equation}
Return the first such vector with
\begin{equation}\label{eq:roundaccept}
 16y_j\ip{s^{(\ell)}}{z_j}\ge n
 \qquad(1\le j\le m).
\end{equation}
If all tables fail, return an arbitrary vector with a failure flag.

\subsection{At most 182 corrections}

Fix $t=16s^*$. The potential $\Phi(a)=\norm{a-t}_2^2$ starts at $256n$.
On a violated observation,
\[
 y_j\ip{t-a}{z_j}>
 (16/5-2)n=6n/5.
\]
Before clipping,
\[
 \Phi(a+y_jz_j)
 =\Phi(a)-2y_j\ip{t-a}{z_j}+\norm{z_j}_2^2
 <\Phi(a)-7n/5.
\]
Clipping coordinatewise to $[-16,16]^n$ cannot increase the distance to $t$,
since $t$ belongs to this box. After $k$ corrections, nonnegativity of the
potential would therefore require $k<1280/7$. Hence at most 182 corrections
are possible, the cap is not reached on promised inputs, and the final
$v=a/16$ satisfies
\begin{equation}\label{eq:roundmean}
 y_j\ip{v}{z_j}\ge n/8
 \qquad(1\le j\le m).
\end{equation}

\subsection{A rounding trial succeeds with probability at least one half}

Condition on the sample and the resulting $a$. For each coordinate,
\[
 \mathbb E s_i^{(\ell)}
 =2\frac{a_i+16}{32}-1=\frac{a_i}{16}.
\]
For a fixed observation, the independent variables
$X_i=y_js_i^{(\ell)}z_{j,i}$ lie in $[-1,1]$ and their sum has mean at least
$n/8$ by \eqref{eq:roundmean}.

We use the classical bounded-variable estimate of Hoeffding~\cite{hoeffding1963}.
For completeness, if a bounded random variable $V$ takes values in an interval
of length $c$, then
\begin{equation}\label{eq:mgf}
 \mathbb E e^{\lambda(V-\mathbb EV)}
 \le e^{\lambda^2c^2/8}.
\end{equation}
Indeed, the second derivative of the logarithm of this moment-generating
function is the variance under an exponentially tilted distribution. A
variable in an interval of length $c$ has variance at most $c^2/4$, as follows
by bounding its squared distance from the midpoint. The logarithm and its first
derivative vanish at zero; integrating the second-derivative bound twice proves
\eqref{eq:mgf}. Multiplication for independent variables and optimized
exponential Markov give
\[
 \Pr\!\left[\sum_i X_i-\mathbb E\sum_i X_i\le-u\right]
 \le e^{-u^2/(2n)}.
\]
Take $u=n/16$. Failure of \eqref{eq:roundaccept} on a fixed observation has
probability at most $e^{-n/512}$. A union bound over the sample gives
\[
 \Pr[\text{one table fails}]
 \le me^{-n/512}\le me^{-b}\le\tfrac12,
\]
because $b=\lceil\log_2(2m)\rceil\ge\ln(2m)$. The tables are independent, so
the probability that all $q$ fail is at most $2^{-q}\le\delta/2$.

\subsection{Generalization from compressed descriptions}

\paragraph{Large-dimensional branch.}
First condition on all the presampled rounding tables, which are independent
of the observations. The corrected vector $a$ can be reconstructed by replaying
the ordered sequence of the at most 182 observations on which corrections
were made. An observation may appear several times in this sequence.
Start at zero, replay the corrections \eqref{eq:clip}, and then apply a
specified rounding table. This fully describes the output.

Fix an ordered sequence $I$ of $k\le182$ sample indices and a table index.
Condition on the observations at the distinct indices in $I$. The reconstructed
classifier is now fixed, and at least $m-k$ remaining observations are still
iid from $\mathcal D$. Therefore
\[
 \Pr[\text{this classifier has error}>\varepsilon
       \text{ and is correct on all }m\text{ observations}]
 \le (1-\varepsilon)^{m-k}
 \le e^{-\varepsilon(m-182)}.
\]
There are at most
\[
 q\sum_{k=0}^{182}m^k\le183q\,m^{182}
\]
descriptions. Every accepted output is one of them and is correct on the
whole sample. A union bound, followed by averaging over the tables, yields
\begin{equation}\label{eq:compressionbound}
 \Pr[\text{accepted output has error}>\varepsilon]
 \le183q\,m^{182}e^{-\varepsilon(m-182)}.
\end{equation}
This argument includes all possible correction-index sequences and all table
indices. It therefore accounts for their adaptive selection; it does not
condition on the event that the actual algorithm chose a particular sequence.

\paragraph{Small-dimensional branch.}
For a fixed sign-vector classifier of error greater than $\varepsilon$, the
probability of correctness on all $m$ observations is at most
$e^{-\varepsilon m}$. There are $2^n$ sign vectors. Since
$n<512\lceil\log_2(2m)\rceil$ and $2^{\lceil\log_2(2m)\rceil}\le4m$,
\begin{equation}\label{eq:smallbound}
 \Pr[\text{accepted output has error}>\varepsilon]
 \le2^n e^{-\varepsilon m}
 \le(4m)^{512}e^{-\varepsilon m}.
\end{equation}

\paragraph{The sample size bounds both probabilities.}
By \eqref{eq:properm}, $m=2^r$ and
$\varepsilon m\ge512(r+2)+q$. Thus
\[
 (4m)^{512}e^{-\varepsilon m}
 =\exp\big(512\ln2\,(r+2)-\varepsilon m\big)
 \le e^{-q}\le\delta/2.
\]
The same sample choice gives $m>q$ and $m>183$. Subtracting the logarithm of
the right side of \eqref{eq:compressionbound} from that of
\eqref{eq:smallbound} gives
\[
 512\ln4+330\ln m-\ln q-\ln183-182\varepsilon
 >
 512\ln4+329\ln m-\ln183-182>0.
\]
Hence the large-branch generalization failure bound is also at most
$\delta/2$. Adding its rounding failure probability at most $\delta/2$ proves
the confidence guarantee. The small branch has no construction failure on
promised inputs.

\subsection{The explicit sample bound and bit complexity}

Let $k=\lceil\log_2(1/\varepsilon)\rceil$ and $A=q+k+12$. Let $m_0=2^{r_0}$
be the least power of two at least $1024A/\varepsilon$. Then
\[
 m_0\le2048A/\varepsilon,\qquad
 r_0\le11+k+\log_2 A.
\]
Since $\log_2 A\le A$ and $q\ge2$,
\[
 2A\ge13+k+\log_2 A+q/512.
\]
For example, replacing $\log_2 A$ by $A$ reduces the last inequality to
$q(511/512)\ge1$. Consequently,
\[
 \varepsilon m_0\ge1024A\ge512(r_0+2)+q.
\]
Minimality in \eqref{eq:properm} gives $m\le m_0$, proving
\eqref{eq:properexplicit} and therefore
\[
 m=O\!\left(
 \frac{1+\log(1/\varepsilon)+\log(1/\delta)}{\varepsilon}
 \right).
\]

If $\varepsilon=u/v$ is rational, the integer $r$ is found by successive exact
comparisons
\[
 u\,2^r\ge v\big(512(r+2)+q\big).
\]
All other logarithmic ceilings are likewise obtained by integer comparisons
to powers of two. The algorithm implements no irrational normalization.

Each $a_i$ is an integer in $[-16,16]$, and a dot product uses
$O(\log(n+1))$ bits. The correction procedure takes at most 183 full scans.
The large branch has a straightforward bit bound
\[
 O\big(nm(183+q)\log((n+1)(m+1))\big).
\]
The small branch checks at most $(4m)^{512}$ candidates, with bit bound
\[
 O\big((4m)^{512}nm\log((n+1)(m+1))\big).
\]
Both are polynomials with absolute constant exponents. Sample and table
storage uses $O(n(m+q))=O(nm)$ bits. The enumeration exponent is large;
this is a polynomial-time existence theorem rather than a practically
attractive replacement for Theorem~\ref{thm:group}.

\paragraph{Optional nonempty output groups.}
If the target is assumed to have both groups nonempty and the output must
also do so, restrict the small-branch enumeration accordingly. In the large
branch, if an accepted vector is constant, flip one fixed coordinate.
This changes each signed training score by at most 2; since
$n\ge512b\ge512$, the original lower bound $n/16$ remains strictly positive.
The repaired output is still correct on the entire sample and is a
deterministic function of the same compressed description. Thus the
generalization argument and sample size are unchanged. Exact 60/40
compatibility of this output remains outside the conclusion.

\section{Exact group-normalized compatibility is NP-complete}
\label{app:hardness}

Consider the following finite decision problem. The input is a list of
$x\in\{0,1\}^n$, without labels. Does some partition $P\sqcup N=[n]$ make every
point satisfy either \eqref{eq:groupplus} or \eqref{eq:groupminus}?

\begin{theorem}\label{thm:npcomplete}
This decision problem is NP-complete. It remains NP-complete when every input
point has exactly $n/2$ active coordinates and there is a connected graph on
the input points whose edges all have Hamming distance exactly 4.
In the reduction below, every YES instance admits a partition for which all
points satisfy the positive branch.
\end{theorem}

The result concerns construction of a compatible partition, not prediction
under the promise. We give the reduction in full.

\subsection{Two consequences of the compatibility inequalities}

For any partition and example, write
\[
 z=2x-\1,\quad
 s_i=\begin{cases}1&i\in P,\\-1&i\in N,\end{cases}
 \quad A=\ip{s}{z},\quad S=\sum_i s_i,\quad Z=\sum_i z_i.
\]
The two inequalities for branch $y\in\{-1,1\}$ are equivalent to
\begin{equation}\label{eq:fullcompat}
 5yA\ge n+|S-5yZ|.
\end{equation}
To verify this directly, use
\[
 |P|=(n+S)/2,\quad |N|=(n-S)/2,
 \quad p=(n+S+Z+A)/4,\quad q=(n-S+Z-A)/4.
\]
For $y=1$, substitution gives
$5A\ge n+S-5Z$ and $5A\ge n-S+5Z$, whose conjunction is
\eqref{eq:fullcompat}. For $y=-1$ it gives
$-5A\ge n+S+5Z$ and $-5A\ge n-S-5Z$, proving the other case.

In particular, $yA\ge n/5$. If $A=n/5$ and $Z=0$ on the positive branch,
then \eqref{eq:fullcompat} forces $S=0$.

A second consequence is that two compatible points on opposite branches
have Hamming distance at least $n/5$. Their positive-group active counts
differ by at least $|P|/5$, and their negative-group counts differ in the
opposite direction by at least $|N|/5$. The sum of these two differences is
bounded above by their Hamming distance.

\subsection{A suitable satisfiability source problem}

Use exact-2-in-4-SAT with signed literals and four distinct variables in each
clause. In sign notation, a clause has the form
\begin{equation}\label{eq:exacttwo}
 \sum_{j=1}^4\lambda_j t_{i_j}=0,
 \qquad t_i\in\{-1,1\},\quad\lambda_j\in\{-1,1\},
\end{equation}
where the $i_j$ are distinct. The equation means that exactly two of the
four signed literals are true.

For completeness, its NP-hardness follows from the standard NP-complete
3-SAT problem~\cite{cook1971}, with three distinct
variables per clause. Introduce a common variable $F$ and replace a clause
$(a\lor b\lor c)$ by $\operatorname{NAE}(a,b,c,F)$, where NAE requires that
the displayed truth values are not all equal. Global complementation of all
variables preserves NAE, including signed literals. A satisfying assignment
may therefore be oriented to make $F$ false; then these requirements are
exactly the original OR clauses.

With a fresh variable $d$ for each such clause, replace
\[
 \operatorname{NAE}(a,b,c,F)
 \quad\text{by}\quad
 \operatorname{NAE}(a,b,d)\ \land\
 \operatorname{NAE}(\neg d,c,F).
\]
If $a=b$, the first condition forces $d$ opposite to them, and the second
condition excludes exactly the case $a=b=c=F$. If $a\ne b$, the first
condition permits either $d$, and a choice satisfying the second always
exists. This proves the equivalence after existentially quantifying $d$.

Finally replace every $\operatorname{NAE}(L_1,L_2,L_3)$ by the exact-two
condition on $(L_1,L_2,L_3,w)$ with a fresh $w$. The sum of the first three
Boolean truth values can be completed to 2 by a Boolean $w$ exactly when
that sum is 1 or 2. All clauses constructed this way have four distinct
variables. Membership in NP is immediate. The usual distinct-variable
version of 3-SAT causes no restriction: a shorter non-tautological clause
can be padded with fresh variables and both choices of their signs.

\subsection{The reduction}

Pad with unused variables until the number $v$ of variables is a multiple
of four and $v\ge12$. Write
\[
 h=v/2,\qquad m_A=v/2,\qquad n=5v/2.
\]
Use $v$ pairs of free coordinates $(i,0),(i,1)$, and $m_A$ anchor coordinates.
Fix an anchor vector $a$ containing $v/4$ plus signs and $v/4$ minus signs.
For every balanced free row $r\in\{-1,1\}^{2v}$ used below, define
\begin{equation}\label{eq:reductionrow}
 z_r=(a,r),\qquad x_r=(\1+z_r)/2.
\end{equation}
Both parts are balanced, so $Z=0$ and every $x_r$ has exactly $n/2$ active
coordinates.

\paragraph{Core rows forcing opposite signs in every free pair.}
For $e\in\{-1,1\}^v$ with $\sum_i e_i=0$, let
\[
 R(e)=(e_1,e_1,e_2,e_2,\ldots,e_v,e_v).
\]
Let $I=\{1,\ldots,h\}$, $J=\{h+1,\ldots,v\}$, and let $e^0$ be plus on $I$
and minus on $J$. Use the following spanning tree of the complete bipartite
graph with sides $I,J$:
\[
 \mathcal T=
 \{(1,j):j\in J\}
 \ \cup\
 \{(i,h+1):i\in I,\ i\ne1\}.
\]
For each edge $(i,j)$, let $e^{ij}$ be obtained from $e^0$ by exchanging its
signs at $i,j$. Put the $v$ rows $R(e^0)$ and $R(e^{ij})$ into the core list,
and put their negatives into the list as well.

\paragraph{Clause rows.}
For a clause \eqref{eq:exacttwo}, set its four participating pairs to
$(\lambda_j,-\lambda_j)$. Among the remaining $v-4$ pairs, use $(1,1)$ on
exactly $(v-4)/2$ pairs and $(-1,-1)$ on the others, in a fixed order.
This gives a balanced row $r_c$. Include both $r_c$ and $-r_c$.

\paragraph{Connecting paths.}
Connect each core row to $R(e^0)$ by adding rows when necessary.
For a clause row or its negative, first replace its four mixed pairs by
two $(1,1)$ pairs and two $(-1,-1)$ pairs. Exactly one coordinate of each
mixed pair changes, so this step has Hamming distance 4 and results in a
balanced pure-pair row $R(e)$.

For a balanced pure-pair row, repeatedly exchange a plus pair outside $I$
with a minus pair inside $I$ until reaching $R(e^0)$. Each exchange changes
four coordinates and preserves balance; at most $h$ exchanges are needed.
All intermediate rows after the first step are pure-pair rows.

If the formula has $c$ clauses, the resulting list contains
$M=O((v+c)v)$ rows. The construction is polynomial. The paths supply a
connected graph on the corresponding points, and every edge has Hamming
distance 4. Crucially,
\[
 4<n/5=v/2
\]
because $v\ge12$.

\subsection{A satisfying assignment gives compatibility}

Let $t$ solve \eqref{eq:exacttwo} for every clause. Set $s=a$ on the anchors,
and set the two free signs of pair $i$ to $(t_i,-t_i)$. Then $S=0$.

For every pure-pair row, the free inner product is zero because each pair
cancels. For a clause row $r_c$, the inactive pure pairs cancel and the
free inner product is
\[
 2\sum_{j=1}^4\lambda_j t_{i_j}=0.
\]
The same holds for its negative. All additional path rows are pure-pair rows,
so the free inner product vanishes on every constructed row. Hence
\[
 \ip{s}{z_r}=m_A=n/5,\qquad Z=0,\qquad S=0
\]
for every point. Equation~\eqref{eq:fullcompat} holds with equality on the
positive branch. More explicitly, both groups have size $n/2$ and
\[
 p(x_r)=3n/10,\qquad q(x_r)=n/5.
\]
Thus every input point is compatible with the positive branch.

\subsection{A compatible partition gives a satisfying assignment}

Suppose one partition makes all constructed points compatible. The branches
are well defined because $n\ge1$. Adjacent points in the path graph cannot
have opposite branches: their Hamming distance is 4, strictly less than $n/5$.
Connectedness forces the same branch everywhere. If it is the negative
branch, interchange $P$ and $N$. We may assume all points use the positive
branch.

Write $s_A,s_F$ for the anchor and free parts. The core list contains
$r^0=R(e^0)$ and $-r^0$. Applying \eqref{eq:fullcompat}, or its consequence
$A\ge n/5=m_A$, gives
\[
 \ip{s_A}{a}+\ip{s_F}{r^0}\ge m_A,\qquad
 \ip{s_A}{a}-\ip{s_F}{r^0}\ge m_A.
\]
Adding shows $\ip{s_A}{a}\ge m_A$. There are only $m_A$ anchor coordinates,
so $s_A=a$ coordinate by coordinate. The two inequalities now force
$\ip{s_F}{r^0}=0$. For this row, $A=m_A=n/5$ and $Z=0$, so
\eqref{eq:fullcompat} implies $S=0$. The anchors are balanced; consequently
$\sum s_F=0$.

Every other core row $r$ also appears with its negative, so the same argument
now gives $\ip{s_F}{r}=0$. Define
\[
 T_i=s_{i,0}+s_{i,1}.
\]
We have $\sum_iT_i=0$. Orthogonality to $R(e^0)$ and $R(e^{ij})$, whose
difference changes the repeated signs at $i,j$, yields $T_i=T_j$ on every
edge of $\mathcal T$. Since $\mathcal T$ is connected, all $T_i$ are equal.
Their sum is zero, so $T_i=0$ for each $i$. Thus
$s_{i,1}=-s_{i,0}$.

Set $t_i=s_{i,0}$. For each clause row, all inactive pure pairs contribute
zero and orthogonality becomes
\[
 2\sum_{j=1}^4\lambda_jt_{i_j}=0.
\]
This gives a satisfying assignment of the exact-2-in-4 instance.

The two directions prove the reduction. A proposed partition is an $n$-bit
certificate, and all its counts and inequalities can be checked in polynomial
time. This proves Theorem~\ref{thm:npcomplete}.

\subsection{Consequence for a learner required to output compatibility}

\begin{corollary}\label{cor:rp}
Suppose a polynomial-time learner, on every perfectly group-compatible
distribution, returns an explicit partition with incompatibility at most
$\varepsilon$ with probability at least $1-\delta$, for arbitrary positive
accuracy and confidence parameters. The learner may also require correctly
labeled examples and may additionally promise small prediction error.
If its runtime is polynomial in $n,1/\varepsilon,\log(1/\delta)$, then
$\NP=\RP$.
\end{corollary}
\begin{proof}
Given the constructed list of $M$ points, pad it by repetitions of its first
point to a list of length
\[
 M'=2^{\lceil\log_2 M\rceil}<2M.
\]
Take the uniform distribution over these $M'$ positions and supply label $+1$
whenever requested. Sampling one position uses exactly $\log_2M'$ fair bits,
so the simulation has a strict polynomial time bound; no rejection sampling
is needed. Every distinct original point has probability at least $1/M'$.
Choose
\[
 \varepsilon=\frac1{2M'},\qquad\delta=\frac14.
\]

In a YES instance, the all-positive labeled distribution is realizable by
the forward direction of the reduction. With probability at least $3/4$ the
learner returns a partition of incompatibility at most $\varepsilon$.
Such a partition must be compatible with every original point, since even
one incompatible point has mass at least $1/M'>\varepsilon$.

Run the learner and verify its returned partition on the entire finite list.
For a NO instance this verification always rejects, regardless of the learner's
behavior. For a YES instance it accepts with probability at least $3/4$.
The chosen inverse accuracy is polynomial in the formula size. If runtime is
only specified on promised inputs, stop the learner at its stated polynomial
clock and reject on timeout. Any sample allowance polynomial in the raw count
benchmark is also polynomial here, since the number of partitions is at most
$2^n$.

We have a randomized polynomial-time algorithm for an NP-complete problem
with no false acceptance and success probability at least $3/4$ on YES inputs.
Thus $\NP\subseteq\RP$. Conversely, a successful polynomial-length random
string is an NP witness for an RP acceptance, giving $\RP\subseteq\NP$.
Therefore $\NP=\RP$.
\end{proof}

The corollary concerns incompatibility of the \emph{output partition}.
It does not obstruct a predictor with small classification error, nor the
signed-literal representation guarantee in Appendix~\ref{app:proper}.
Indeed, all examples in the YES direction of this reduction have the same
label, while finding a compatible partition still encodes the satisfiability
problem.

\begin{thebibliography}{99}
\bibitem{bb2007}
Avrim Blum and Maria-Florina Balcan.
\emph{Open Problems in Efficient Semi-Supervised PAC Learning}.
COLT 2007, Sections 3 and 3.2.
\url{https://home.ttic.edu/~avrim/Papers/openproblems.pdf}.

\bibitem{bb2010}
Maria-Florina Balcan and Avrim Blum.
\emph{A Discriminative Model for Semi-Supervised Learning}.
Journal of the ACM 57(3), Article 19, pp.~1--46, 2010.
DOI: \href{https://doi.org/10.1145/1706591.1706599}{10.1145/1706591.1706599}.
Author manuscript, 2009, Sections 2, 3.1.1, and 6.2.
\url{https://home.ttic.edu/~avrim/Papers/labeled_unlabeled-2009.pdf}.

\bibitem{bbel2013}
Maria-Florina Balcan, Christopher Berlind, Steven Ehrlich, and Yingyu Liang.
\emph{Efficient Semi-supervised and Active Learning of Disjunctions}.
ICML 2013, Proceedings of Machine Learning Research 28(1), pp.~633--641.
\url{https://proceedings.mlr.press/v28/balcan13.pdf}.

\bibitem{bbel2013supp}
Maria-Florina Balcan, Christopher Berlind, Steven Ehrlich, and Yingyu Liang.
Supplementary material to \cite{bbel2013}, Section 2, Theorem 4.
\url{https://proceedings.mlr.press/v28/balcan13-supp.pdf}.

\bibitem{bb2014}
Maria-Florina Balcan and Christopher Berlind.
\emph{A New Perspective on Learning Linear Separators with Large
$L_qL_p$ Margins}.
AISTATS 2014, Proceedings of Machine Learning Research 33, pp.~68--76.
\url{https://proceedings.mlr.press/v33/balcan14.pdf}.

\bibitem{bb2014supp}
Maria-Florina Balcan and Christopher Berlind.
Supplementary material to \cite{bb2014}, Appendix A.2.
\url{https://proceedings.mlr.press/v33/balcan14-supp.pdf}.

\bibitem{novikoff1962}
A. B. J. Novikoff.
\emph{On Convergence Proofs for Perceptrons}.
Presented at the Symposium on the Mathematical Theory of Automata,
April 1962; pp.~615--622 in the archived manuscript.
\url{https://cs.uwaterloo.ca/~y328yu/classics/novikoff.pdf}.

\bibitem{fw1995}
Sally Floyd and Manfred Warmuth.
\emph{Sample Compression, Learnability, and the Vapnik--Chervonenkis Dimension}.
Machine Learning 21, pp.~269--304, 1995.
\url{https://doi.org/10.1007/BF00993593}.

\bibitem{hoeffding1963}
Wassily Hoeffding.
\emph{Probability Inequalities for Sums of Bounded Random Variables}.
Journal of the American Statistical Association 58(301), pp.~13--30, 1963.
\url{https://doi.org/10.1080/01621459.1963.10500830}.

\bibitem{cook1971}
Stephen A. Cook.
\emph{The Complexity of Theorem-Proving Procedures}.
Proceedings of the Third Annual ACM Symposium on Theory of Computing,
pp.~151--158, 1971.
\url{https://doi.org/10.1145/800157.805047}.

\bibitem{lmw2002}
Ming Li, Bin Ma, and Lusheng Wang.
\emph{On the Closest String and Substring Problems}.
Journal of the ACM 49(2), pp.~157--171, 2002.
\url{https://doi.org/10.1145/506147.506150}.
\end{thebibliography}

\end{document}
